\documentclass[11pt]{article}
\usepackage[margin=1in]{geometry}
\usepackage{amsmath,amssymb}
\usepackage[hidelinks]{hyperref}

\title{Literature Answer: Parasol Graphs Do Not Characterize Rolewicz Beta}
\author{Run \texttt{fa\_banach\_001}}
\date{2026-06-15}

\newcommand{\dist}{\operatorname{dist}}

\begin{document}
\maketitle

\paragraph{Status.}
\texttt{literature\_already\_answered}.  The extracted problem from
arXiv:1504.04250 has a negative answer in later literature.

\paragraph{Source question.}
In Baudier--Zhang, \emph{\((\beta)\)-distortion of some infinite graphs},
arXiv:1504.04250, Section ``Metric characterization of asymptotic properties'',
the authors prove that if a Banach space \(Y\) admits an equivalent norm with
Rolewicz property \((\beta)\), then
\[
  \sup_{\ell\ge 1} c_Y(P^\omega_\ell)=\infty .
\]
They then ask whether the converse holds:
\[
  Y\notin\langle(\beta)\rangle
  \quad\Longrightarrow\quad
  \sup_{\ell\ge 1} c_Y(P^\omega_\ell)<\infty .
\]

\paragraph{Later answer.}
Baudier--Causey--Dilworth--Kutzarova--Randrianarivony--Schlumprecht--Zhang,
\emph{On the geometry of the countably branching diamond graphs},
arXiv:1612.01984, Section 5 defines a family \(\mathcal G^\omega\) of recursive
finite-height bundle graphs and explicitly includes the countably branching
parasol graphs \(P^\omega_k\).  Their Theorem 5.1 states that if a Banach
space \(Y\) admits an equivalent asymptotically midpoint uniformly convex
norm, then
\[
  \sup_{k\in\mathbb N} c_Y(G^\omega_k)=\infty
  \qquad ((G^\omega_k)\in\mathcal G^\omega).
\]
Immediately after the proof, the authors point out that this theorem gives a
negative answer to Problem 4.1 of Baudier--Zhang.

\paragraph{Concrete witness.}
The classical James space \(\mathcal J\) is nonreflexive, so it cannot admit
an equivalent norm with property \((\beta)\), because property \((\beta)\)
implies reflexivity.  On the other hand, \(\mathcal J\) has a
2-asymptotically uniformly convex norm; AUC implies AMUC.  Applying Theorem 5.1
to the parasol sequence gives
\[
  \sup_{k\ge 1} c_{\mathcal J}(P^\omega_k)=\infty .
\]
Thus \(\mathcal J\notin\langle(\beta)\rangle\) while the parasol distortions are
unbounded, so the converse asked in arXiv:1504.04250 is false.

\paragraph{Why this is not new.}
The supporting paper explicitly identifies the answer: it names the
Baudier--Zhang problem and states that the theorem gives a negative answer.
The James-space facts are included only to make the counterexample mechanism
auditable.

\paragraph{References.}
\begin{itemize}
\item F. P. Baudier and S. Zhang, \emph{\((\beta)\)-distortion of some infinite graphs}, arXiv:1504.04250.
\item F. P. Baudier, R. Causey, S. Dilworth, D. Kutzarova, N. L. Randrianarivony, T. Schlumprecht, and S. Zhang, \emph{On the geometry of the countably branching diamond graphs}, arXiv:1612.01984.
\item F. P. Baudier and G. Lancien, \emph{Geometric nonlinear functional analysis}, arXiv:2512.00817, Appendix on James spaces.
\end{itemize}

\end{document}
